%一些常用的宏定义
\newcommand{\bbc}{{\mathbb{C}}}
\newcommand{\bbr}{\mathbb{R}}
\newcommand{\bbq}{\mathbb{Q}}
\newcommand{\bbz}{\mathbb{Z}}
\newcommand{\bbn}{\mathbb{N}}
\newcommand{\bbd}{\mathbb{D}}
\newcommand{\bbh}{\mathbb{H}}
\newcommand{\bba}{\mathbb{A}}
\newcommand{\bbp}{\mathbb{P}}
\newcommand{\bbf}{\mathbb{F}}


\newcommand{\ten}{\otimes}
\newcommand{\Var}{\mathbf{Var}}


\newcommand{\calo}{\mathcal{O}}
\newcommand{\cali}{\mathcal{I}}


\newcommand{\fraka}{{\mathfrak{a}}}
\newcommand{\frakm}{{\mathfrak{m}}}
\newcommand{\frakp}{{\mathfrak{p}}}


\newcommand{\Frac}{\mathrm{Frac}}
\newcommand{\Der}{\operatorname{Der}}
\newcommand{\Spec}{\operatorname{Spec}}
\newcommand{\spec}{\operatorname{Spec}}
\newcommand{\mSpec}{\operatorname{mSpec}}
\newcommand{\depth}{\operatorname{depth}}
\newcommand{\idealht}{\operatorname{ht}}
\newcommand{\codim}{\operatorname{codim}}
\newcommand{\Supp}{\operatorname{Supp}}
\newcommand{\trdeg}{\operatorname{trdeg}}
\newcommand{\Ass}{\operatorname{Ass}}
\newcommand{\Ann}{\operatorname{Ann}}


\newcommand{\kk}{\mathsf{k}}
\newcommand{\kkk}{\mathbf{k}}
\newcommand{\KK}{\mathsf{K}}
\newcommand{\KKK}{\mathbf{K}}
\newcommand{\rkk}{\kappa} % residue field
\newcommand{\fkk}{\mathscr{K}} % fraction field
\renewcommand{\d}{\mathrm{d}}
\renewcommand{\i}{\mathrm{i}}
\renewcommand{\P}{\partial}


% \renewcommand{\ker}{\mathrm{Ker}\ }
% \newcommand{\ord}{\mathrm{ord}}
% \renewcommand{\hom}{\mathrm{Hom}}
% \renewcommand{\gcd}{\mathrm{g.c.d.}}
% \newcommand{\laplace}{\Delta}
% \newcommand{\lcm}{\mathrm{l.c.m.}}
\newcommand{\mat}[4]{\left( \begin{array}{cc}#1 &#2 \\ #3 &#4\end{array}\right)}
% \renewcommand{\vec}[1]{\boldsymbol{#1}}
% \renewcommand{\proofname}{\indent\it Proof}




\newcommand{\contradiction}{
    \raisebox{-0.6ex}{\makebox[2.4ex][c]{\begin{tikzpicture}
        \draw (0,0) circle (1.2ex);
        \draw (0.7 ex, 0.7 ex) -- (-0.4 ex, 0.7 ex);
        \draw (-0.4 ex, 0.7 ex) -- (0.7 ex, -0.4 ex);
        \draw (0.7 ex, -0.4 ex) -- (0.7 ex, 0.7 ex);
        \draw (0.7 ex, 0.7 ex) -- (-0.848 ex, -0.848 ex);
    \end{tikzpicture}}}
    \ \ 
}

% legendre符号
\makeatletter
\def\legendre@dash#1#2{\hb@xt@#1{%
  \kern-#2\p@
  \cleaders\hbox{\kern.5\p@
    \vrule\@height.2\p@\@depth.2\p@\@width\p@
    \kern.5\p@}\hfil
  \kern-#2\p@
  }}
\def\@legendre#1#2#3#4#5{\mathopen{}\left(
  \sbox\z@{$\genfrac{}{}{0pt}{#1}{#3#4}{#3#5}$}%
  \dimen@=\wd\z@
  \kern-\p@\vcenter{\box0}\kern-\dimen@\vcenter{\legendre@dash\dimen@{#2}}\kern-\p@
  \right)\mathclose{}}
\newcommand\legendre[2]{\mathchoice
  {\@legendre{0}{1}{}{#1}{#2}}
  {\@legendre{1}{.5}{\vphantom{1}}{#1}{#2}}
  {\@legendre{2}{0}{\vphantom{1}}{#1}{#2}}
  {\@legendre{3}{0}{\vphantom{1}}{#1}{#2}}
}
\def\dlegendre{\@legendre{0}{1}{}}
\def\tlegendre{\@legendre{1}{0.5}{\vphantom{1}}}
\makeatother